\documentclass[12pt]{article}
\usepackage[utf8]{inputenc}
\usepackage[T1]{fontenc}
\usepackage{hyperref}
\title{Zapotec, Zaachila Language Resources for K-12 Curriculum Development}
\author{Compiled by Claude Sonnet 5.5 (claude-sonnet-5-5) for Andr\'as Kornai}
\date{}
\begin{document}
\maketitle

\section{Introduction}
Zaachila Zapotec (ISO 639-3: ztx; Glottocode zaac1239) is an endangered Otomanguean language of Oaxaca, Mexico, spoken as a first language mainly by older adults. It is not known to be taught in schools. The searches found very little material specific to this variety. This report lists only what turned up and does not pad the list.

\section{Primary Educational Resources}
No primers, textbooks or school programmes for Zaachila Zapotec were found in search results.

\section{Linguistic and Reference Materials}
\begin{itemize}
\item Wikipedia article on Zaachila Zapotec (English overview of the variety): \url{https://en.wikipedia.org/wiki/Zaachila_Zapotec}
\item Dictionary of the Zapotec Language, hosted by AILLA (University of Texas). It is a multi-volume Zapotec--Spanish and Spanish--Zapotec dictionary drawing on many communities. It is not specific to Zaachila, and I did not confirm that it includes Zaachila forms: \url{https://islandora-ailla.lib.utexas.edu/islandora/object/ailla%3A279600}
\end{itemize}

\section{Religious and Community Resources}
\begin{itemize}
\item Global Recordings Network page for Zapotec, Zaachila \& Jalpan, which offers Bible-teaching audio and a slideshow: \url{https://globalrecordings.net/en/language/ztx}
\end{itemize}

\section{Digital and Audio Resources}
The Global Recordings Network audio listed above is the only audio found. No corpus, software or video resources specific to this variety were found.

\section{Curriculum Development Context}
Very little documented material exists for Zaachila Zapotec, and the speakers are mostly elderly. Any K-12 curriculum would therefore need new documentation work with remaining speakers first. Related Zapotec varieties have more material, but they should not be treated as interchangeable with Zaachila.
\end{document}